\section{Low-dimensional target representations with arbitrarily deep suffixes}
\label{sec:critical-low-rank-targets}

A low-rank first layer gives another target class for which the
critical learner has a polynomial sample bound without a restriction
on the prior. The teacher's remaining depth and width can be arbitrary.
The useful parameters are its projection rank and its product norm.
The theorem below follows the variational route and has a sufficiently
large width qualification. Corollary~\ref{cor:critical-all-width-structured-targets}
gives its every-realizable-width version using the same target compiler.

\begin{theorem}[Low-rank target representations]
\label{thm:critical-low-rank-targets}
Suppose a Boolean target $f$ has a finite tanh representation with
$f(x)z(x)\ge1$ on the full cube. Let $P\ge1$ bound the product of
its layer Frobenius norms, and suppose its first weight matrix,
including the constant input coordinate, has rank at most $r\ge1$.
Put
\[
 S=1+\log(d+L+2)+r\log(P+r+2).
\]
There are universal constants $C_0,C_1,C_2$ such that
\[
 L\ge C_0S\quad\Longrightarrow\quad
 J_L(f)\le V_*+C_1S^2/L.
\]
Under that depth condition, a sufficient sample size for the original
critical posterior is
\[
 n\ge\frac{C_2}{\varepsilon}
 \left\{(d+L+2)^{C_2}(P+r+2)^{C_2r}
                    +\log(1/\delta)\right\}.
\]
The posterior conclusion and width-first quantifiers are those of
Theorem~\ref{thm:critical-excess-learning}. In particular, the width
threshold is finite and independent of the realized sample, but is
not quantified here. For fixed $r$, the bound is polynomial and
independent of the teacher's depth and width.
\end{theorem}

\begin{proof}
Constant functions have constant-size threshold implementations, so
consider a nonconstant target. There are at most $P^2$ relevant
coordinates. Indeed, if coordinate $i$ is relevant, a bichromatic
cube edge and the network's Lipschitz bound give
\[
 \|W^1e_i\|_2\prod_{j\ge2}\|W^j\|_F\ge1.
\]
Square and sum this inequality over relevant coordinates.

Freeze every irrelevant coordinate to an independent uniform sign.
Its first-layer column is absorbed into the bias column. The expected
squared Frobenius norm of the resulting first matrix is the original
squared norm, so some freezing does not increase it. The new matrix's
range is contained in the original range, and therefore its rank
also does not increase. The function is unchanged. Write $k\le P^2$
for the reduced input dimension, $m=k+1$ and $R=\sqrt m$.
We keep the original ambient $d$ in the learner's prior normalization.

Factor the reduced first matrix as $UV^{\mathsf T}$, where the columns
of $V$ are orthonormal and their number is at most $r$; decrease $r$
for the construction if necessary. The network output has the form
\[
 z(x)=g(V^{\mathsf T}\bar x),\qquad \bar x=(x,1),
\]
where $g$ is $P$-Lipschitz. This follows because
$\|U\|_{\rm op}=\|W^1\|_{\rm op}$ and every tanh map is
one-Lipschitz. Oppositely labeled inputs consequently obey
\begin{equation}\label{eq:low-rank-projected-separation}
 \|V^{\mathsf T}\bar x-V^{\mathsf T}\bar y\|_2\ge2/P.
\end{equation}

Choose powers of two $a,\eta$ with
\[
 \frac1{128P\sqrt r}<a\le\frac1{64P\sqrt r},\qquad
 \frac{a}{200m}<\eta\le\frac{a}{100m}.
\]
Round each entry of $V^{\mathsf T}$ to a multiple of $\eta$, obtaining
$B$. On every cube input this changes each projection coordinate by
at most $m\eta/2\le a/200$. Each coordinate of $B\bar x$ is an
integer multiple of $\eta$, and $a/\eta$ is an integer.

For each positive input choose a nearest center in the grid
$a\mathbb Z^r$, with arbitrary deterministic tie breaking. Retain
all centers obtained in this way. Their number satisfies
\[
 M\le(CPR\sqrt r)^r\le(CP^2\sqrt r)^r.
\]
For each selected center $c$, define a box predicate by
$|(B\bar x)_j-c_j|\le a$ for every $j$. Every positive input lies
in its selected box. No negative input lies in a selected box:
if $y$ is a positive input selecting $c$ and $x$ lies in that box,
their rounded projections differ in every coordinate by at most
$3a/2$. Their true projections then differ in Euclidean norm by less
than $2a\sqrt r\le1/(32P)$, contradicting
\eqref{eq:low-rank-projected-separation}.
Thus the target is the disjunction of these box predicates.

The construction has bounded integer coefficients and internal
margins. On the projection lattice, the lower and upper box
inequalities have the same signs as
\[
 (B\bar x)_j-(c_j-a-\eta/2),\qquad
 (c_j+a+\eta/2)-(B\bar x)_j,
\]
respectively. Every score has absolute value at least $\eta/2$.
Multiplying by $2/\eta$ gives integer coefficients and margin at
least one. Their norms are bounded by a universal polynomial in
$P,r$, since $m\le P^2+1$ and $|c_j|\le R+a$.
An AND of the $2r$ signed inequalities is implemented by the score
$\sum_i t_i-(2r-1)$, and an OR of the $M$ signed box outputs by
$\sum_c u_c+(M-1)$. These also have integer margin at least one.
Constant rails supply the offsets.

We have obtained a threshold circuit with three nonconstant stages,
width at most $C(rM+P^2+1)$ and coefficient norms bounded by a fixed
polynomial in $P,r,M$. Apply the controlled circuit seed theorem
and the steering-after-prefix lemma. The logarithmic resource
parameter is at most
\[
 C\{1+\log(d+L+2)+\log(M+2)+\log(P+r+2)\}\le CS.
\]
The original input coordinates can be restored as zero columns;
only their Gaussian normalization contributes the displayed $d$.
This proves the rate and learner-depth bounds. Finally,
Theorem~\ref{thm:critical-excess-learning} bounds the sample factor
by a polynomial prefactor times $e^{CS}$, which has the form stated
in the theorem after increasing $C_2$.
\end{proof}

\begin{corollary}[An approximate first-layer bottleneck]
\label{cor:critical-approx-low-rank-targets}
The theorem also applies when the first matrix $W^1$ has a rank-$r$
truncated singular-value decomposition $A$ such that
\[
 \sqrt{d+1}\,\|W^1-A\|_{\rm op}
              \prod_{j\ge2}\|W^j\|_{\rm op}\le1/8,
\]
with $P$ replaced by $8P/7$.
\end{corollary}

\begin{proof}
The replacement changes every output by at most $1/8$ and does not
increase the first Frobenius norm. Multiply the readout by $8/7$
to restore unit margin, and apply the theorem.
\end{proof}

This is a restriction on a target representation, not on prior
support. It does not show that product norm alone permits a critical
rate excess of order $\log^2(d+L+P)/L$. A product bound controls the
sum of logarithms of layer stable ranks, but a margin-preserving
rank-one truncation generally needs a much smaller discarded
singular value than that budget alone supplies.
